\scanpage{5}
1°) \uline{$n$ pair}$=2\nu$,\quad \uline{$n-1$ impair}
\qquad $n-1=2\nu-1=2(\nu-1)+1$
\diagram{\[
\left.\begin{array}{ll}
H^*(X_0):&
1,\xi_0,\ldots,\xi_0^{\nu-1},
[\xi_0^\nu,\lambda'_\nu],\frac12\xi_0^{\nu+1},\ldots,\frac12\xi_0^{2\nu-1}\\[5pt]
H^*(\bar X_1):&
1,\xi_1,\ldots,\xi_1^{\nu-1},
\frac12\xi_1^\nu,\frac12\unsure{\xi_1^{\nu+1}},\ldots,\frac12\xi_1^{2\nu-1}
\end{array}\right]
\quad
\begin{array}{l}\xi_0\rightsquigarrow\xi_1\\\lambda'_\nu\rightsquigarrow0\end{array}
\]}
Donc\quad $H^i(X_0)\mathbin{-}H^i(\bar X_1)$
\uline{bijectif} \deleted{excepté} si $i\ne2\nu=n$

$H^{2\nu}(X_0)\mathbin{-}H^{2\nu}(\bar X_1)$
\uline{surjectif}, noyau engendré par $\lambda'_\nu$.
\[
\text{Donc}\qquad
a)\left\{\begin{array}{ll}
\Phi^i=0&\deleted{\text{sauf}}\ \text{si }i\ne(2\nu=n)\\
\Phi^{\unsure{2\nu|n}}\simeq\unsure{\mu|M}(-\nu)&
\text{(isom. déterminé au signe près\ldots)}
\end{array}\right.
\]
\[
b)\left\{\begin{array}{l}
R^if_*\text{ constant si }i\ne n\\
R^nf_*\text{ a un sous-faisceau concentré en o, isom. à }\unsure{\mu|M}(\ldots),\\
\text{le quotient par ce dernier est constant}
\end{array}\right.
\]
\margin{c) $\pi$ opère trivialement sur $H^*(\bar X_1)$ et sur $\Phi^{2\nu}$}

\figure{\begin{tikzpicture}
\draw (0,0)--(8,.3) .. controls (8.3,.3) and (8.1,-.4) .. (8.4,-.5)--(10.4,-.5);
\end{tikzpicture}}
2°) $n$ impair $=2\nu+1$,\qquad $n-1=2\nu$
\diagram{\[
\left.\begin{array}{ll}
H^*(X_0):&
1,\xi_0,\ \xi_0^{\nu-1},\ \xi_0^\nu,\ \frac12\xi_0^{\nu+1},
\ldots,\frac12\xi_0^{2\nu}\\[5pt]
H^*(\bar X_1):&
1,\xi_1,\ldots,\xi_1^{\nu-1},[\xi_1^\nu,\lambda_\nu],
\frac12\xi_1^{\nu+1},\ldots,\frac12\xi_1^{2\nu}
\end{array}\right]
\quad\xi_0\rightsquigarrow\xi_1
\]}
Donc\quad $H^i(X_0)\mathbin{-}H^i(\bar X_1)$ bijectif si $i\ne2\nu=n-1$

$H^{2\nu}(X_0)\mathbin{-}H^{2\nu}(\bar X_1)$ injectif,
conoyau engendré par $\deleted{\unsure{\xi}}\inserted{\varphi_\nu}$.
\[
a)\left\{\begin{array}{ll}
\Phi^i=0&\text{si }i\ne2\nu+1=n\\
\Phi^n\simeq\unsure{\mu|M}(\unsure{?}\ldots)&
\end{array}\right.
\]
\[
b)\left\{\begin{array}{l}
R^if_*\text{ constant si }i\ne n-1\\
R^{n-1}f_*\text{ est ``pur''}\ldots
\end{array}\right.
\]
\margin{c) $\pi$ opère trivialement sur les $H^i(\bar X_1)$ pour $i\ne n-1$,
mais non sur $\Phi^n$ \ldots}
\margin{4}
